\section{Заключение}\label{sec:conclusion}
Для минимизации суммы выпуклых $(L_0,L_1)$-гладких функций метод YETI
достигает верхней оценки $O(T_0\log^2N)$, а для любого корректного метода
доказана нижняя оценка $\Omega(T_0/\log^2N_0)$. Если $\Delta/\varepsilon\le N_0^c$
при некотором фиксированном $c\ge1$, то эти оценки совпадают с точностью
до степеней $\log N_0$. Без ограничения на рост начальной невязки $\Delta$
оптимальность верхней оценки не установлена.

На 65 задачах при относительной точности $10^{-4}$ испытанный вариант YETI
занял первое место (в том числе разделив его с другими методами)
в 46 случаях, а его критерий остановки по истории запросов сработал
на всех 65 задачах. Основной выигрыш в этих опытах даёт эвристика
(приложение~\ref{app:experiments}). Однако её матричные операции могут
быть дорогими, поэтому выигрыш в числе запросов не обязательно означает
такой же выигрыш во времени.

\subsection{Открытые вопросы}\label{sec:open-questions}
\begin{enumerate}
\item \emph{Большая начальная невязка.}
Какова оптимальная зависимость сложности от $\Delta$?
Если $\log(\Delta/\varepsilon)$ имеет порядок $L_1R$, то полученные верхняя
и нижняя оценки могут различаться полиномиально по $L_1R$. Нужно либо
улучшить верхнюю оценку, либо доказать нижнюю оценку, учитывающую $\Delta$.
При умеренных $\Delta$ остаётся устранить логарифмический разрыв.

\item \emph{Индивидуальные константы гладкости.}
Можно ли получить более точную оценку сложности для заданного набора
$(L_{0,i},L_{1,i})$, не заменяя все $L_{1,i}$ их максимумом?
Неравномерный выбор компонент уже позволяет обойтись средним
значением $L_{0,i}$ (приложение~\ref{app:variants}). Чтобы улучшить общую
оценку, нужно учесть связь между $L_{1,i}$ и невязками отдельных слагаемых
и построить соответствующую нижнюю оценку: наша нижняя оценка получена
на подклассе с одинаковыми $L_{0,i}$.

\item \emph{Ограничения и неевклидова геометрия.}
Как перенести результат на минимизацию по выпуклому множеству
и на методы с расстоянием Брэгмана? Зеркальный спуск при обобщённой
гладкости изучен в работе~\citet{yu2026mirror}. Здесь дополнительно
потребуются анализ уменьшения дисперсии, замена евклидова потенциала
и оценка числа отсечений.

\item \emph{Негладкие слагаемые.}
Можно ли получить аналогичные двусторонние оценки для конечных сумм,
у которых нормы субградиентов ограничены через невязку?
Такая негладкая стохастическая постановка рассматривается в работе
\citet{kovalev2026nonsmooth}. Для конечной суммы нужно отдельно определить
модель оракула и выяснить, какой выигрыш даёт уменьшение дисперсии.

\item \emph{Методы второго и более высоких порядков.}
Можно ли совместить уменьшение дисперсии для конечных сумм
\citep{tsedenov2026second} с ускорением при обобщённой гладкости
гессиана \citep{agafonov2026inexact}? Для этого нужны оценки ошибок
стохастических производных и новый локальный решатель. Условия на старшие
производные не следуют из предположения~\ref{asm:smoothness}.

\item \emph{Двойственная задача.}
Можно ли интерпретировать YETI как покоординатный метод для двойственной
задачи и связать его со схемой из работы \citet{podlipnova2026sparse}?
Помимо преобразования задачи, потребуется сопоставить доступные оракулы
и стоимость одного шага: запрос компоненты $f_i$ и запрос координаты
двойственного градиента в общем случае не эквивалентны.
\end{enumerate}

Отдельно остаётся найти сложность вне области действия
предположения~\ref{asm:accuracy}, в частности при $\delta=0$.
Этот вопрос касается неоднородности слагаемых и не совпадает
с вопросом о большой начальной невязке $\Delta$.
